\scanpage{290}
\mspage{31}

On a le diagramme de foncteurs canoniques
\diagram{\begin{tikzcd}[ampersand replacement=\&,column sep=large]
\underline{\mathrm{Ty}} \arrow[r,hook] \arrow[d,hook] \& {}_{\varepsilon}\backslash\mathrm{Cat} \arrow[d]\\
\underline{\mathrm{Typ}} \arrow[r,arrows={Hooks[harpoon]-}] \& \mathrm{Cat}\mathpunct{,}
\end{tikzcd}}
où les flèches horizontales sont pl. fidèles.

5. \uline{Propriétés} \deleted{de $\underline{\mathrm{Ty}}$, \unsure{et}} \uline{des morphismes d'intervalles}

Je me suis plus ou moins convaincu que pour se prêter à des calculs et à des raisonnements simples, il faut encore plus restreindre les morphismes dans la \inserted{définition de la} catégorie $\underline{\mathrm{Ty}}$. Je \deleted{\unsure{vais}} exige\deleted{\unsure{r}} que les morphismes $\tau\to\tau'$ soient des foncteurs
\[
\unsure{w|u}\colon I_\tau\mathrel{\unsure{\text{---}|\longrightarrow}}I_{\tau'},
\]
qui sont compatibles avec la structure d'\uline{intervalles} \uline{orientés} de $I_\tau,I_{\tau'}$, i.e. \inserted{qu'ils} transforment source en source, but en but, mais de plus qu'ils soient \uline{croissants} sur les objets, (pour la relation d'ordre de progression des chemins).
C'est une restriction qui géométriquement me parait assez naturelle. (Elle est similaire à la restriction analogue, dans la définition pré-catégorique des morphismes entre $\mathbb{\Delta}_{\unsure{i|1}}\ldots$) C'est cette hypothèse qui permet de faire une théorie de structure assez simple des \deleted{foncteurs} morphismes d'intervalles. Mais pour l'instant je vais distinguer ``morphismes'' et ``morphismes croissants''
% archive sheet number: 289
